% HAFF: Holographic Alaya-Field Framework
% Book-Level Abstract for Collected Volume
% Standalone Document

\documentclass[12pt,a4paper]{article}
\usepackage{amsmath,amssymb,amsfonts,amsthm}
\usepackage{physics}
\usepackage{hyperref}
\usepackage{geometry}
\geometry{margin=1in}

\title{The Abstract of the Holographic Alaya–Field Framework}

\author{
  Sidong Liu, PhD \\
  iBioStratix Ltd \\
  \texttt{sidongliu@hotmail.com}
}

\date{July 2026}

\begin{document}

\maketitle

% ============================================================================
% ABSTRACT
% ============================================================================

\section*{Abstract}

This volume develops a structural framework for understanding the emergence of spacetime, locality, and classicality within quantum theory, without assuming any preferred subsystem decomposition or fundamental spacetime background.

The framework proposes a programme-level way to study features traditionally treated as primitive---such as spatial separation, temporal ordering, and observer--system distinctions---through coarse-grained, physically accessible observable subalgebras. Building on insights from algebraic quantum theory, quantum information, and holography \cite{Haag1996,Zanardi2001,Zanardi2004,RyuTakayanagi2006,VanRaamsdonk2010}, it treats observable algebras as a useful structural focus, not as a completed derivation of all spacetime or observational structure from one principle.

Within this approach, subsystems are not assumed as canonical tensor factors but are modelled through stable patterns of accessibility, constrained by interaction structure, information-theoretic robustness, and environmental redundancy \cite{Zurek2009}. Geometry, time, and measurement are analyzed in regimes where such structures persist, and the resulting descriptions are contingent, emergent, and structurally bounded within the stated assumptions. Different coarse-graining schemes can give rise to inequivalent effective geometries once the relevant reconstruction data are supplied, even when defined on the same underlying quantum state, without modifying the underlying dynamics.

The framework does not propose new fundamental laws or cosmological dynamics. Instead, it reorganizes explanatory priority, clarifying some conditions under which physical descriptions become available. By identifying structural limits of locality, temporality, and observation within the programme, this work aims to mark where the present framework stops and where further dynamics, dictionaries, or external assumptions would be needed.

\bigskip

\noindent\textbf{Keywords}: emergent geometry, accessible algebras, coarse-graining, quantum foundations, algebraic quantum theory, structural realism

\newpage

% ============================================================================
% REFERENCES
% ============================================================================

\begin{thebibliography}{99}

\bibitem{Haag1996}
R. Haag, \emph{Local Quantum Physics: Fields, Particles, Algebras}, Springer-Verlag (1996).

\bibitem{Zanardi2001}
P. Zanardi, \emph{Virtual Quantum Subsystems}, Phys. Rev. Lett. \textbf{87}, 077901 (2001).

\bibitem{Zanardi2004}
P. Zanardi, D. A. Lidar, and S. Lloyd, \emph{Quantum Tensor Product Structures are Observable Induced}, Phys. Rev. Lett. \textbf{92}, 060402 (2004).

\bibitem{RyuTakayanagi2006}
S. Ryu and T. Takayanagi, \emph{Holographic Derivation of Entanglement Entropy from AdS/CFT}, Phys. Rev. Lett. \textbf{96}, 181602 (2006).

\bibitem{VanRaamsdonk2010}
M. Van Raamsdonk, \emph{Building up spacetime with quantum entanglement}, Gen. Relativ. Gravit. \textbf{42}, 2323 (2010).

\bibitem{Zurek2009}
W. H. Zurek, \emph{Quantum Darwinism}, Nature Physics \textbf{5}, 181 (2009).

\end{thebibliography}

\end{document}
